\ifdefined\gkver\else\def\gkver{student}\fi
\documentclass[\gkver]{gaokaozhenti}
\begin{document}

\chapter{2012年海南}

\begin{enumerate}
\item
%% number: 1
%% paperName: 2012年海南高考第 1 题 3 分
%% typeId: 1
%% type: 单选题
%% chapter: 运动和力
%% point: 牛顿第二定律
%% method: 无
%% score: 3
%% degree: 650
%% duplicateId: 0
%% body:
根据牛顿第二定律，下列叙述正确的是\xzanswer{D}
\fourchoices[answer=D]
{物体加速度的大小跟它的质量和速度大小的乘积成反比}
{物体所受合力必须达到一定值时，才能使物体产生加速度}
{物体加速度的大小跟它所受作用力中的任一个的大小成正比}
{当物体质量改变但其所受合力的水平分力不变时，物体水平加速度大小与其质量成反比}
%% answer: D
%% memo:
\memoanswer{
物体加速度的大小与物体受到的合力成正比，与物体的质量成反比，选项A错误；力是产生加速度的原因，只要有合力，物体就有加速度，它们之间有瞬时对应关系，不存在累积效应，选项B错误；物体加速度的大小与它受到的合力成正比，选项C错误；根据矢量的合成和分解得$F_{x}=ma_{x}$，当$F_{x}$一定时，物体水平加速度的大小与其质量成反比，选项D正确。
}

\item
%% number: 2
%% paperName: 2012年海南高考第 2 题 3 分
%% typeId: 1
%% type: 单选题
%% chapter: 磁场
%% point: 带电粒子在复合场中的运动
%% method: 无
%% score: 3
%% degree: 650
%% duplicateId: 0
%% body:
如图所示，在两水平极板间存在匀强电场和匀强磁场，电场方向竖直向下，磁场方向垂直于纸面向里，一带电粒子以某一速度沿水平直线通过两极板，若不计重力，下列四个物理量中哪一个改变时，粒子运动轨迹不会改变\xzanswer{B}
\onepicture{figs/02.png}
\fourchoices[answer=B]
{粒子速度的大小}
{粒子所带的电荷量}
{电场强度}
{磁感应强度}
%% answer: B
%% memo:
\memoanswer{
粒子能沿水平直线通过两极板，说明粒子所受的电场力与洛伦兹力平衡，有$qE=qvB$。当仅改变粒子所带的电荷量时，等式两边同时随$q$变化，平衡仍然成立，粒子的运动轨迹不发生改变；而改变粒子速度的大小、电场强度或磁感应强度，都会破坏这一平衡，使粒子的运动轨迹发生改变。故正确选项为B。
}

\item
%% number: 3
%% paperName: 2012年海南高考第 3 题 3 分
%% typeId: 1
%% type: 单选题
%% chapter: 电场
%% point: 带电粒子的直线运动
%% method: 无
%% score: 3
%% degree: 650
%% duplicateId: 0
%% body:
如图所示，直线上有$o$、$a$、$b$、$c$四点，$ab$间的距离与$bc$间的距离相等，在$o$点处有固定点电荷，已知$b$点电势高于$c$点电势。若一带电负电荷的粒子仅在电场力作用下先从$c$点运动到$b$点，再从$b$点运动到$a$点，则\xzanswer{C}
\onepicture{figs/03.png}
\fourchoices[answer=C]
{两过程中电场力做的功相等}
{前一过程中电场力做的功大于后一过程中电场力做的功}
{前一过程中，粒子电势能不断减小}
{后一过程中，粒子动能不断减小}
%% answer: C
%% memo:
\memoanswer{
根据点电荷电场的电场线及电势分布可知，$b$点电势高于$c$点，则$O$点固定的是正电荷。负电荷移动过程中电场力做功$W=qU$，又$U_{bc}<U_{ab}$，则前一过程中电场力做功小于后一过程中电场力做功，A、B选项错误；两个过程中电场力均做正功，电势能减小、动能增加，C正确，D错误。正确选项为C。
}

\item
%% number: 4
%% paperName: 2012年海南高考第 4 题 3 分
%% typeId: 1
%% type: 单选题
%% chapter: 交流电
%% point: 变压器
%% method: 无
%% score: 3
%% degree: 650
%% duplicateId: 0
%% body:
如图所示，理想变压器原、副线圈匝数比为$20:1$，两个标有“$12\UV$，$6\UW$”的小灯泡并联在副线圈的两端，当两灯泡都正常工作时，原线圈电路中电压表和电流表（可视为理想的）的示数分别是\xzanswer{D}
\onepicture{figs/04.png}
\fourchoices[answer=D]
{$120\UV$，$0.10\UA$}
{$240\UV$，$0.025\UA$}
{$120\UV$，$0.05\UA$}
{$240\UV$，$0.05\UA$}
%% answer: D
%% memo:
\memoanswer{
根据电压关系$\frac{U_{1}}{U_{2}}=\frac{n_{1}}{n_{2}}$，由$U_{2}=12\UV$可知$U_{1}=240\UV$；根据功率关系$P_{1}=P_{2}$，$P_{1}=U_{1}I_{1}$，$P_{2}=12\UW$，联立可得$I_{1}=0.05\UA$。正确选项为D。
}

\item
%% number: 5
%% paperName: 2012年海南高考第 5 题 3 分
%% typeId: 1
%% type: 单选题
%% chapter: 电磁感应
%% point: 楞次定律推广
%% method: 无
%% score: 3
%% degree: 650
%% duplicateId: 0
%% body:
如图所示，一质量为$m$的条形磁铁用细线悬挂在天花板上，细线一水平金属圆环中穿过。现将环从位置Ⅰ释放，环经过磁铁到达位置Ⅱ。设环经过磁铁上端附近时细线的张力分别为$T_{1}$和$T_{2}$，重力加速度大小为$g$，则\xzanswer{A}
\onepicture{figs/05.png}
\fourchoices[answer=A]
{$T_{1}>mg$，$T_{2}>mg$}
{$T_{1}<mg$，$T_{2}<mg$}
{$T_{1}>mg$，$T_{2}<mg$}
{$T_{1}<mg$，$T_{2}>mg$}
%% answer: A
%% memo:
\memoanswer{
由楞次定律可知，圆环下落通过磁铁的过程中，穿过圆环的磁通量变化，圆环中产生感应电流，圆环与磁铁间存在相互作用的安培力。在Ⅰ位置，圆环受向上的安培力，磁铁受向下的安培力，故细线张力$T_{1}>mg$；在Ⅱ位置，圆环同样受向上的安培力，磁铁受向下的安培力，细线张力$T_{2}>mg$。故正确选项为A。
}

\item
%% number: 6
%% paperName: 2012年海南高考第 6 题 3 分
%% typeId: 1
%% type: 单选题
%% chapter: 运动和力
%% point: 牛顿定律中的图象
%% method: 无
%% score: 3
%% degree: 650
%% duplicateId: 0
%% body:
如图所示，表面处处同样粗糙的楔形木块$abc$固定在水平地面上，$ab$面和$bc$面与地面的夹角分别为$\alpha$和$\beta$，且$\alpha>\beta$，一初速度为$v_{0}$的小物块沿斜面$ab$向上运动，经时间$t_{0}$后到达顶点$b$时，速度刚好为零；然后让小物块立即从静止开始沿斜面$bc$下滑。在小物块从$a$运动到$c$的过程中，可能正确描述其速度大小$v$与时间$t$的关系的图像是\xzanswer{C}
\onepicture[h=2.6cm, align=right]{figs/06.png}
\fourchoices[answer=C, ispicture=true, h=2.2cm]
{figs/06a.png}
{figs/06b.png}
{figs/06c.png}
{figs/06d.png}
%% answer: C
%% memo:
\memoanswer{
由于阻力做功，物块下滑到地面时速度$v$小于初始速度$v_{0}$，物块在斜面上做匀变速运动，但前后两段受力不同，加速度不同，平均速度不同：上升阶段平均速度$\overline{v_{1}}=\frac{v_{0}}{2}$，下降阶段平均速度$\overline{v_{2}}=\frac{v}{2}$，则$\overline{v_{1}}>\overline{v_{2}}$；且$\alpha>\beta$，$\sin\alpha>\sin\beta$，$x_{ab}<x_{bc}$，上滑时间$t_{1}=\frac{x_{ab}}{\overline{v_{1}}}$，下滑时间$t_{2}=\frac{x_{bc}}{\overline{v_{2}}}$，$t_{1}<t_{2}$。综上可判断正确选项为C。
}

\item
%% number: 7
%% paperName: 2012年海南高考第 7 题 4 分
%% typeId: 2
%% type: 多选题
%% chapter: 功和能
%% point: 动能定理
%% method: 无
%% score: 4
%% degree: 650
%% duplicateId: 0
%% body:
下列关于功和机械能的说法，正确的是\xzanswer{BC}
\fourchoices[answer=BC]
{在有阻力作用的情况下，物体重力势能的减少不等于重力物体所做的功}
{合力对物体所做的功等于物体动能的改变量}
{物体的重力势能是物体与地球之间的相互作用能，其大小与势能零点的选取有关}
{运动物体动能的减少量一定等于其重力势能的增加量}
%% answer: BC
%% memo:
\memoanswer{
重力势能的减少量等于重力对物体所做的功，与有无阻力作用无关，A错误；由动能定理可知，合力对物体所做的功等于物体动能的变化量，B正确；物体的重力势能是物体与地球之间的相互作用能，其大小与势能零点的选取有关，C正确；只有在只有重力做功的前提下，物体动能的减少量才等于其重力势能的增加量，D错误。正确选项为BC。
}

\item
%% number: 8
%% paperName: 2012年海南高考第 8 题 4 分
%% typeId: 2
%% type: 多选题
%% chapter: 力物体的平衡
%% point: 摩擦力
%% method: 无
%% score: 4
%% degree: 650
%% duplicateId: 0
%% body:
下列关于摩擦力的说法，正确的是\xzanswer{CD}
\fourchoices[answer=CD]
{作用在物体上的滑动摩擦力只能使物体减速，不可能使物体加速}
{作用在物体上的静摩擦力只能使物体加速，不可能使物体减速}
{作用在物体上的滑动摩擦力既可能使物体减速，也可能使物体加速}
{作用在物体上的静摩擦力既可能使物体加速，也可能使物体减速}
%% answer: CD
%% memo:
\memoanswer{
摩擦力的方向可能与物体运动方向相同，也可能与物体运动方向相反，但一定是与物体相对运动方向或相对运动趋势方向相反。当滑动摩擦力的方向与物体的运动方向相同时，使物体加速，选项A错误、C正确；当静摩擦力的方向与物体的运动方向相反时，使物体减速，选项B错误、D正确。
}

\item
%% number: 9
%% paperName: 2012年海南高考第 9 题 4 分
%% typeId: 2
%% type: 多选题
%% chapter: 电路
%% point: 电容
%% method: 无
%% score: 4
%% degree: 650
%% duplicateId: 0
%% body:
将平行板电容器两极板之间的距离、电压、电场强度大小和极板所带的电荷分别用$d$、$U$、$E$和$Q$表示．下列说法正确的是\xzanswer{AD}
\fourchoices[answer=AD]
{保持$U$不变，将$d$变为原来的两倍，则$E$变为原来的一半}
{保持$E$不变，将$d$变为原来的一半，则$U$变为原来的两倍}
{保持$d$不变，将$Q$变为原来的两倍，则$U$变为原来的一半}
{保持$d$不变，将$Q$变为原来的一半，则$E$变为原来的一半}
%% answer: AD
%% memo:
\memoanswer{
由$E=\frac{U}{d}$，保持$U$不变，将$d$变为原来的两倍，$E$变为原来的一半，A正确；保持$E$不变，将$d$变为原来的一半，则$U$变为原来的一半，B错误；由$C=\frac{Q}{U}$，$C=\frac{\varepsilon S}{4\pi kd}$，保持$d$不变，$C$不变，$Q$加倍，$U$加倍，C错误；又$E=\frac{U}{d}=\frac{Q}{Cd}=\frac{4\pi kQ}{\varepsilon S}$，则$Q$变为原来的一半，$E$变为原来的一半，D正确。正确选项为AD。
}

\item
%% number: 10
%% paperName: 2012年海南高考第 10 题 4 分
%% typeId: 2
%% type: 多选题
%% chapter: 磁场
%% point: 左手定则
%% method: 无
%% score: 4
%% degree: 650
%% duplicateId: 0
%% body:
图中装置可演示磁场对通电导线的作用，电磁铁上下两磁极之间某一水平面内固定两条平行金属导轨，$L$是置于导轨上并与导轨垂直的金属杆，当电磁铁线圈两端$a$、$b$，导轨两端$e$、$f$，分别接到两个不同的直流电源上时，$L$便在导轨上滑动，下列说法正确的是\xzanswer{BD}
\onepicture{figs/10.png}
\fourchoices[answer=BD]
{若$a$接正极，$b$接负极，$e$接正极，$f$接负极，则$L$向右滑动}
{若$a$接正极，$b$接负极，$e$接负极，$f$接正极，则$L$向右滑动}
{若$a$接负极，$b$接正极，$e$接正极，$f$接负极，则$L$向左滑动}
{若$a$接负极，$b$接正极，$e$接负极，$f$接正极，则$L$向左滑动}
%% answer: BD
%% memo:
\memoanswer{
若$a$接正极，$b$接负极，则磁场向上，$e$接正极、$f$接负极时导体棒中电流向外，根据左手定则，导体棒受力向左，向左滑动，A错误；若$a$接正极，$b$接负极，磁场向上，$e$接负极、$f$接正极，导体棒中电流向里，根据左手定则，导体棒受力向右，向右滑动，B正确；若$a$接负极，$b$接正极，磁场向下，$e$接正极、$f$接负极，导体棒中电流向外，根据左手定则，导体棒受力向右，向右滑动，C错误；若$a$接负极，$b$接正极，磁场向下，$e$接负极、$f$接正极，导体棒中电流向里，根据左手定则，导体棒受力向左，向左滑动，D正确。正确选项为BD。
}

\item
%% number: 11
%% paperName: 2012年海南高考第 11 题 4 分
%% typeId: 3
%% type: 填空题
%% chapter: 曲线运动
%% point: 人造地球卫星
%% method: 无
%% score: 4
%% degree: 650
%% duplicateId: 0
%% body:
地球同步卫星到地心的距离$r$可用地球质量$M$、地球自转周期$T$与引力常量$G$表示为$r=$\_\_\_\_\_\_\_\_\_\_\_\_\_。
%% answer: 
\jdanswer{
$\sqrt[3]{\frac{GMT^{2}}{4\pi^{2}}}$
}
%% memo:
\memoanswer{
提供同步卫星做圆周运动的向心力是万有引力，根据牛顿第二定律得$G\frac{Mm}{r^{2}}=m\left(\frac{2\pi}{T}\right)^{2}r$，所以$r=\sqrt[3]{\frac{GMT^{2}}{4\pi^{2}}}$。
}

\item
%% number: 12
%% paperName: 2012年海南高考第 12 题 4 分
%% typeId: 3
%% type: 填空题
%% chapter: 电场
%% point: 电场强度
%% method: 对称法
%% score: 4
%% degree: 450
%% duplicateId: 0
%% body:
$N$（$N>1$）个电荷量均为$q$（$q>0$）的小球，均匀分布在半径为$R$的圆周上，示意如图，若移去位于圆周上$P$点的一个小球，则圆心$O$点处的电场强度大小为\_\_\_\_\_\_\_\_，方向\_\_\_\_\_\_\_\_\_\_。（已知静电力常量为$k$）
\onepicture[align=right]{figs/12.png}
%% answer: 
\jdanswer{
$\frac{kq}{R^{2}}$，沿$OP$指向$P$点
}
%% memo:
\memoanswer{
移去$P$点小球后，余下小球在$O$点产生的合场强与$P$点小球在$O$点产生的场强等大反向，故大小为$\frac{kq}{R^{2}}$，方向沿$OP$指向$P$点。
}

\item
%% number: 13
%% paperName: 2012年海南高考第 13 题 6 分
%% typeId: 4
%% type: 实验题
%% chapter: 电路
%% point: 电路实验
%% method: 无
%% score: 6
%% degree: 650
%% duplicateId: 0
%% body:
图示电路可用来测量电阻的阻值。其中$E$为电源，$R$为已知电阻，$R_{x}$为待测电阻，$V$可视为理想电压表，$\mathrm{S}_{0}$为单刀单掷开关，$\mathrm{S}_{1}$、$\mathrm{S}_{2}$为单刀双掷开关。
\begin{enumerate}
	\item
	当$\mathrm{S}_{0}$闭合时，若$\mathrm{S}_{1}$、$\mathrm{S}_{2}$均向左闭合，电压表读数为$U_{1}$；若$\mathrm{S}_{1}$、$\mathrm{S}_{2}$均向右闭合，电压读数为$U_{2}$。由此可求出$R_{x}=$\_\_\_\_\_\_\_\_\_。
	\item
	若电源电动势$E=1.5\UV$，内阻可忽略；电压表量程为$1\UV$，$R=100\UO$。此电路可测量的$R_{x}$的最大值为\_\_\_\_\_\_\_$\UO$。
\end{enumerate}
\onepicture[align=right]{figs/13.png}
%% answer: 
\jdanswer{
\begin{enumerate}
	\item
	$\frac{U_{1}}{U_{2}}R$

	\item
	200

\end{enumerate}
}
%% memo:
\memoanswer{
$\mathrm{S}_{1}$、$\mathrm{S}_{2}$均向左闭合时$U_{1}=IR_{x}$，$\mathrm{S}_{1}$、$\mathrm{S}_{2}$均向右闭合时$U_{2}=IR$，故$R_{x}=\frac{U_{1}}{U_{2}}R$；电源电动势为$1.5\UV$，电压表量程为$1\UV$，则$\frac{U_{1}}{U_{2}}\leq\frac{2}{1}$，故$R_{x}$的最大值为$200\UO$。
}

\item
%% number: 14
%% paperName: 2012年海南高考第 14 题 9 分
%% typeId: 4
%% type: 实验题
%% chapter: 功和能
%% point: 动能势能
%% method: 无
%% score: 9
%% degree: 450
%% duplicateId: 0
%% body:
一水平放置的轻弹簧，一端固定，另一端与一小滑块接触，但不粘连；初始时滑块静止于水平气垫导轨上的$O$点，如图（a）所示。现利用此装置探究弹簧的弹性势能$E_{p}$与其被压缩时长度的改变量$x$的关系，先推动小滑块压缩弹簧，用米尺测出$x$的数值；然后将小滑块从静止释放。用计时器测出小滑块从$O$点运动至气垫导轨上另一固定点A所用的时间$t$。多次改变$x$，测得的$x$值及其对应的$t$值如下表所示。（表中的$1/t$值是根据$t$值计算得出的）。
\onepicture[h=2.2cm, align=right]{figs/14a.png}

\begin{center}
\begin{tabular}{|c|c|c|c|c|c|}
\hline
$x$（$\Ucm$） & 1.00 & 1.50 & 2.00 & 2.50 & 3.00 \\ \hline
$t$（$\Us$） & 3.33 & 2.20 & 1.60 & 1.32 & 1.08 \\ \hline
$1/t$（$\Us^{-1}$） & 0.300 & 0.455 & 0.625 & 0.758 & 0.926 \\ \hline
\end{tabular}
\end{center}

（1）根据表中数据，在图（b）中的方格纸上作图线。
\onepicture[h=3.8cm, align=right]{figs/14b.png}

（2）回答下列问题：（不要求写出计算或推导过程）

①已知点（0，0）在$(1/t)\text{-}x$图线上，从$(1/t)\text{-}x$图线看，$1/t$与$x$是什么关系？

②从理论上分析，小滑块刚脱离弹簧时的动能$E_{k}$与$1/t$是什么关系（不考虑摩擦力）？

③当弹簧长度改变量为$x$时，弹性势能$E_{p}$与相应的$E_{k}$是什么关系？

④综合考虑以上分析，$E_{p}$与$x$是什么关系？
%% answer: 
\jdanswer{
\begin{enumerate}
	\item
	$(1/t)\text{-}x$图线为过原点的直线，如图所示。

	\item
	①$1/t$与$x$成正比；②$E_{k}$与$(1/t)^{2}$成正比；③$E_{p}=E_{k}$；④$E_{p}$与$x^{2}$成正比。

\end{enumerate}
}
%% memo:
\memoanswer{
$(1/t)$与$x$的关系图线为过原点的直线，故$1/t$与$x$成正比。由$v=\frac{x}{t}$，动能$E_{k}=\frac{1}{2}mv^{2}=\frac{1}{2}m\left(\frac{x}{t}\right)^{2}$，故$E_{k}\propto\left(\frac{1}{t}\right)^{2}$。根据机械能守恒，弹簧形变量为$x$时对应的弹性势能与弹簧弹出时物块的动能相等，即$E_{p}=E_{k}$。综上可得$E_{p}\propto x^{2}$。
\onepicture[align=right]{figs/14_an.png}
}

\item
%% number: 15
%% paperName: 2012年海南高考第 15 题 9 分
%% typeId: 6
%% type: 计算题
%% chapter: 曲线运动
%% point: 竖直平面圆周运动
%% method: 无
%% score: 9
%% degree: 450
%% duplicateId: 0
%% body:
如图，在竖直平面内有一固定光滑轨道，其中$AB$是长为$R$的水平直轨道，$BCD$是圆心为$O$、半径为$R$的$3/4$圆弧轨道，两轨道相切于$B$点。在外力作用下，一小球从$A$点由静止开始做匀加速直线运动，到达$B$点时撤除外力。已知小球刚好能沿圆轨道经过最高点$C$，重力加速度大小为$g$。求：
\begin{enumerate}
	\item
	小球在$AB$段运动的加速度的大小；
	\item
	小球从$D$点运动到$A$点所用的时间。
\end{enumerate}
\onepicture[align=right]{figs/15.png}
%% answer: 
\jdanswer{
\begin{enumerate}
	\item
	$\frac{5}{2}g$

	\item
	$(\sqrt{5}-\sqrt{3})\sqrt{\frac{R}{g}}$

\end{enumerate}
}
%% memo:
\memoanswer{
（1）小球在$BCD$段运动时，受到重力$mg$、轨道正压力$N$的作用，如图所示。据题意，$N\geq0$，且小球在最高点$C$所受轨道正压力为零
$N_{C}=0$ ①
设小球在$C$点的速度大小为$v_{C}$，根据牛顿第二定律有
$mg=m\frac{v_{C}^{2}}{R}$ ②
小球从$B$点运动到$C$点，机械能守恒。设$B$点处小球的速度大小为$v_{B}$，有
$\frac{1}{2}mv_{B}^{2}=\frac{1}{2}mv_{C}^{2}+mg\cdot2R$ ③
由于小球在$AB$段由静止开始做匀加速运动，设加速度大小为$a$，由运动学公式有
$v_{B}^{2}=2aR$ ④
由②③④式得：$a=\frac{5}{2}g$ ⑤

（2）设小球在$D$处的速度大小为$v_{D}$，下落到$A$点时的速度大小为$v$，由机械能守恒有
$\frac{1}{2}mv_{B}^{2}=\frac{1}{2}mv_{D}^{2}+mgR$ ⑥
$\frac{1}{2}mv_{B}^{2}=\frac{1}{2}mv^{2}$ ⑦
设从$D$点运动到$A$点所用的时间为$t$，由运动学公式得
$gt=v-v_{D}$ ⑧
由④⑤⑥⑦⑧式得$t=(\sqrt{5}-\sqrt{3})\sqrt{\frac{R}{g}}$ ⑨
\onepicture[align=right]{figs/15_an.png}
}

\item
%% number: 16
%% paperName: 2012年海南高考第 16 题 10 分
%% typeId: 6
%% type: 计算题
%% chapter: 磁场
%% point: 带电粒子在磁场中的运动
%% method: 无
%% score: 10
%% degree: 250
%% duplicateId: 0
%% body:
图（a）所示的$xOy$平面处于匀强磁场中，磁场方向与$xOy$平面（纸面）垂直，磁感应强度$B$随时间$t$变化的周期为$T$，变化图线如图（b）所示。当$B$为$+B_{0}$时，磁感应强度方向指向纸外。在坐标原点$O$有一带正电的粒子$P$，其电荷量与质量之比恰好等于$2\pi/(TB_{0})$。不计重力。设$P$在某时刻$t_{0}$以某一初速度沿$y$轴正向自$O$点开始运动，将它经过时间$T$到达的点记为$A$。
\begin{enumerate}
	\item
	若$t_{0}=0$，则直线$OA$与$x$轴的夹角是多少？
	\item
	若$t_{0}=T/4$，则直线$OA$与$x$轴的夹角时多少？
	\item
	为了使直线$OA$与$x$轴的夹角为$\pi/4$，在$0<t_{0}<T/4$的范围内，$t_{0}$应取何值？
\end{enumerate}
\twopicture[showlabel=false, align=right]
{figs/16a.png}
{figs/16b.png}
%% answer: 
\jdanswer{
\begin{enumerate}
	\item
	0

	\item
	$\pi/2$

	\item
	$T/8$

\end{enumerate}
}
%% memo:
\memoanswer{
（1）设粒子$P$的质量、电荷量与初速度分别为$m$、$q$与$v$，粒子$P$在洛伦兹力作用下在$xOy$平面内做圆周运动，分别用$R$与$T'$表示圆周的半径和运动周期，则有
$qvB=m\left(\frac{2\pi}{T'}\right)^{2}R$ ①
$v=\frac{2\pi R}{T'}$ ②
由①②式与已知条件得$T'=T$ ③
粒子$P$在$t=0$到$t=T/2$时间内沿顺时针方向运动半个圆周，到达$x$轴上$B$点，此时磁场方向反转；继而，在$t=T/2$到$t=T$时间内沿逆时针方向运动半个圆周，到达$x$轴上$A$点，如图（a）所示。$OA$与$x$轴的夹角
$\theta=0$ ④

（2）粒子$P$在$t_{0}=T/4$时刻开始运动，在$t=T/4$到$t=T/2$时间内沿顺时针方向运动$1/4$个圆周，到达$C$点，此时磁场方向反转；继而，在$t=T/2$到$t=T$时间内沿逆时针方向运动半个圆周，到达$B$点，此时磁场方向再次反转；在$t=T$到$t=5T/4$时间内沿顺时针方向运动$1/4$个圆周，到达$A$点，如图（b）所示。由几何关系可知，$A$点在$y$轴上，即$OA$与$x$轴的夹角
$\theta=\frac{\pi}{2}$ ⑤

（3）若在任意时刻$t=t_{0}$（$0<t_{0}<T/4$）粒子$P$开始运动，在$t=t_{0}$到$t=T/2$时间内沿顺时针方向做圆周运动到达$C$点，圆心$O'$位于$x$轴上，圆弧$\widehat{OC}$对应的圆心角为
$\angle OO'C=\frac{2\pi}{T}\left(\frac{T}{2}-t_{0}\right)$ ⑥
此时磁场方向反转；继而，在$t=T/2$到$t=T$时间内沿逆时针方向运动半个圆周，到达$B$点，此时磁场方向再次反转；在$t=T$到$t=T+t_{0}$时间内沿顺时针方向做圆周运动到达$A$点，设圆$O''$，圆弧$\widehat{BA}$对应的圆心角为
$\angle BO''A=\frac{2\pi}{T}t_{0}$ ⑦
如图（c）所示，由几何关系可知，$C$、$B$均在$O'O''$连线上，且$OA\parallel O'O''$ ⑧
若要$OA$与$x$轴成$\pi/4$角，则有$\angle OO'C=\frac{3\pi}{4}$ ⑨
联立⑥⑨式可得
$t_{0}=\frac{T}{8}$ ⑩
\threepicture[showlabel=false, align=right, h=2.2cm]
{figs/16_an1.png}
{figs/16_an2.png}
{figs/16_an3.png}
}

\item
%% number: 17
%% paperName: 2012年海南高考第 17 题 4 分
%% typeId: 2
%% type: 多选题
%% chapter: 气体性质
%% point: 分子力
%% method: 无
%% score: 4
%% degree: 650
%% duplicateId: 0
%% body:
两分子间的斥力和引力的合力$F$与分子间距离$r$的关系如图中曲线所示，曲线与$r$轴交点的横坐标为$r_{0}$。相距很远的两分子在分子力作用下，由静止开始相互接近。若两分子相距无穷远时分子势能为零，下列说法正确的是\_\_\_\_\_\_\_\_。
\onepicture{figs/17.png}
\fivechoices[answer=ACE]
{在$r>r_{0}$阶段，$F$做正功，分子动能增加，势能减小}
{在$r<r_{0}$阶段，$F$做负功，分子动能减小，势能也减小}
{在$r=r_{0}$时，分子势能最小，动能最大}
{在$r=r_{0}$时，分子势能为零}
{分子动能和势能之和在整个过程中不变}
%% answer: ACE
%% memo:
\memoanswer{
$r>r_{0}$时分子力表现为引力，做正功，分子势能减小、动能增加；$r=r_{0}$时分子势能最小、动能最大；$r<r_{0}$时分子力表现为斥力，做负功，分子势能增加、动能减小；分子动能和势能之和在整个过程中保持不变，故A、C、E正确。正确选项为ACE。
}

\item
%% number: 17
%% paperName: 2012年海南高考第 17 题 8 分
%% typeId: 6
%% type: 计算题
%% chapter: 气体性质
%% point: 玻意耳定律
%% method: 无
%% score: 8
%% degree: 450
%% duplicateId: 0
%% body:
如图所示，一气缸水平固定在静止的小车上，一质量为$m$、面积为$S$的活塞将一定量的气体封闭在气缸内，平衡时活塞与气缸底相距$L$。现让小车以一较小的水平恒定加速度向右运动，稳定时发现活塞相对于气缸移动了距离$d$。已知大气压强为$P_{0}$，不计气缸和活塞间的摩擦；且小车运动时，大气对活塞的压强仍可视为$P_{0}$；整个过程温度保持不变。求小车加速度的大小。
\onepicture[align=right]{figs/17b.png}
%% answer: 
\jdanswer{
$\frac{p_{0}Sd}{m(L-d)}$
}
%% memo:
\memoanswer{
设小车加速度大小为$a$，稳定时气缸内气体的压强为$P_{1}$，活塞受到气缸内外气体的压力分别为
$F_{1}=P_{1}S$ ①
$F_{0}=P_{0}S$ ②
由牛顿第二定律得$F_{1}-F_{0}=ma$ ③
小车静止时，在平衡情况下，气缸内气体的压强应为$P_{0}$。由玻意耳定律得
$P_{1}V_{1}=P_{0}V$ ④
式中
$V=SL$ ⑤
$V_{1}=S(L-d)$ ⑥
联立①②③④⑤⑥式得$a=\frac{p_{0}Sd}{m(L-d)}$ ⑦
}

\item
%% number: 18
%% paperName: 2012年海南高考第 18 题 4 分
%% typeId: 3
%% type: 填空题
%% chapter: 机械波
%% point: 波长频率波速
%% method: 无
%% score: 4
%% degree: 650
%% duplicateId: 0
%% body:
某波源$S$发出一列简谐横波，波源$S$的振动图像如图所示。在波的传播方向上有$A$、$B$两点，他们到$S$的距离分别为$45\Um$和$55\Um$。测得$A$、$B$两点开始振动的时间间隔为$1.0\Us$。由此可知：

①波长$\lambda=$\_\_\_\_\_\_\_\_\_$\Um$；

②当$B$点离开平衡位置的位移为$+6\Ucm$时，$A$点离开平衡位置的位移是\_\_\_\_\_\_\_$\Ucm$。
\onepicture[align=right]{figs/18.png}
%% answer: 
\jdanswer{
20，$-6$
}
%% memo:
\memoanswer{
由题意可知波速$v=\frac{x}{t}=10\Ums$，周期$T=2.0\Us$，则波长$\lambda=vT=20\Um$。$A$、$B$两点相距$10\Um$，恰好为半个波长，两点的振动情况完全相反，故$B$点位移为$+6\Ucm$时，$A$点位移为$-6\Ucm$。
}

\item
%% number: 18
%% paperName: 2012年海南高考第 18 题 8 分
%% typeId: 6
%% type: 计算题
%% chapter: 光学
%% point: 全反射
%% method: 无
%% score: 8
%% degree: 650
%% duplicateId: 0
%% body:
一玻璃三棱镜，其横截面为等腰三角形，顶角$\theta$为锐角，折射率为$\sqrt{2}$。现在横截面内有一光线从其左侧面上半部分射入棱镜。不考虑棱镜内部的反射。若保持入射线在过入射点的法线的下方一侧（如图所示），且要求入射角为任何值的光线都会从棱镜的右侧面射出，则顶角$\theta$可在什么范围内取值？
\onepicture[align=right]{figs/18b.png}
%% answer: 
\jdanswer{
$0^\circ<\theta<45^\circ$
}
%% memo:
\memoanswer{
设入射光线经玻璃折射时，入射角为$i$，折射角为$\gamma$，射到棱镜右侧面的入射角为$\alpha$。根据折射定律有
$\sin i=n\sin\gamma$ ①
由几何关系得$\theta=\alpha+\gamma$ ②
当$i=0$时，由①式知$\gamma=0$，$\alpha$有最大值$\alpha_{m}$（如图），由②式得
$\theta=\alpha_{m}$ ③
同时$\alpha_{m}$应小于玻璃对空气的全反射临界角，即
$\sin\alpha_{m}<\frac{1}{n}$ ④
由①②③④式和题给条件可得，棱镜顶角$\theta$的取值范围为
$0^\circ<\theta<45^\circ$ ⑤
\onepicture[align=right]{figs/18_an.png}
}

\item
%% number: 19
%% paperName: 2012年海南高考第 19 题 4 分
%% typeId: 2
%% type: 多选题
%% chapter: 光学
%% point: 光电效应
%% method: 无
%% score: 4
%% degree: 650
%% duplicateId: 0
%% body:
产生光电效应时，关于逸出光电子的最大初动能$E_{k}$，下列说法正确的是\_\_\_\_\_\_。
\fivechoices[answer=ADE]
{对于同种金属，$E_{k}$与照射光的强度无关}
{对于同种金属，$E_{k}$与照射光的波长成反比}
{对于同种金属，$E_{k}$与光照射的时间成正比}
{对于同种金属，$E_{k}$与照射光的频率成线性关系}
{对于不同种金属，若照射光频率不变，$E_{k}$与金属的逸出功成线性关系}
%% answer: ADE
%% memo:
\memoanswer{
由光电方程$E_{k}=h\nu-W=h\frac{c}{\lambda}-W$可知，同种金属的逸出功相同，最大初动能与照射光的强度无关，A正确；最大初动能与照射光的频率成线性关系，与波长不是反比关系，B错误；最大初动能与照射光的时间无关，C错误；对于同种金属，$E_{k}$与照射光的频率成线性关系，D正确；对于不同种金属，保持入射光频率不变，最大初动能$E_{k}$与逸出功成线性关系，E正确。正确选项为ADE。
}

\item
%% number: 19
%% paperName: 2012年海南高考第 19 题 8 分
%% typeId: 6
%% type: 计算题
%% chapter: 运动和力
%% point: 动量守恒定律
%% method: 无
%% score: 8
%% degree: 450
%% duplicateId: 0
%% body:
一静止的$\ce{^{238}_{92}U}$核经$\alpha$衰变成为$\ce{^{234}_{90}Th}$核，释放出的总动能为$4.27\UMeV$。问此衰变后$\ce{^{234}_{90}Th}$核的动能为多少$\UMeV$（保留1位有效数字）？
%% answer: 
\jdanswer{
$0.07\UMeV$
}
%% memo:
\memoanswer{
据题意知，此$\alpha$衰变的衰变方程为$\ce{^{238}_{92}U -> ^{234}_{90}Th + ^{4}_{2}He}$，根据动量守恒定律得
$m_{\alpha}v_{\alpha}=m_{Th}v_{Th}$ ①
式中，$m_{\alpha}$和$m_{Th}$分别为$\alpha$粒子和$\mathrm{Th}$核的质量，$v_{\alpha}$和$v_{Th}$分别为$\alpha$粒子和$\mathrm{Th}$核的速度大小。由题设条件知$\frac{1}{2}m_{\alpha}v_{\alpha}^{2}+\frac{1}{2}m_{Th}v_{Th}^{2}=E_{k}$ ②
$\frac{m_{\alpha}}{m_{Th}}=\frac{4}{234}$ ③
由①②③式得：$\frac{1}{2}m_{Th}v_{Th}^{2}=\frac{m_{\alpha}}{m_{\alpha}+m_{Th}}E_{k}$ ④
代入数据得，衰变后$\ce{^{234}_{90}Th}$核的动能$\frac{1}{2}m_{Th}v_{Th}^{2}=0.07\UMeV$ ⑤
}

\end{enumerate}

\end{document}
